\section*{Source notes and attribution}
\addcontentsline{toc}{section}{Source notes and attribution}
This is an edited condensation of the v139 continuation manuscript and
its earlier notebook, with the same pinned native BTBMS operation.
Historical progress reports, superseded restricted cases and bounded
experimental logs are omitted. The principal arguments are included
in the text; no numbered notebook section is a necessary external
proof reference. The accompanying source concordance records the
origin of the consolidated arguments and the input hashes.

The three prior amendments remain explicit: the actual work model
retains its required elementary relations; the empty-tail-prefix
interval uses the entry count $\mu$; and a proper generated restriction
assumes source Gen. Low coding of actual incoming graphs is shortened
using that same source Gen and the independently proved finite-term
transport. Its scope is the actual legal maps used in this proof;
it is not a claim that an arbitrary weak cut has a canonical graph.

Standard background comprises constructibility, small-structure
satisfaction, admissible Collection, condensation and the
Silver-indiscernible shift obtained from $0^\#$. The following sources
are background references, not published proofs or endorsements of
the BTBMS argument:
\begin{enumerate}[label={[\arabic*]}]
\item Klaus Sutner, \emph{The Lattice of $\beta$ Recursively Enumerable
Sets}, Section 2.2. \url{https://www.cs.cmu.edu/~sutner/papers/beta.pdf}
\item Victoria Gitman, Joel David Hamkins and Thomas A. Johnstone,
\emph{What is the theory ZFC without power set?}
\url{https://victoriagitman.github.io/files/ZFC-.pdf}
\item Victoria Gitman and P. D. Welch, \emph{Ramsey-like Cardinals II}.
\url{https://arxiv.org/abs/1104.4448}
\item Noam Greenberg, \emph{Two Applications of Admissible Computability}.
\url{https://homepages.ecs.vuw.ac.nz/~greenberg/Papers/72-Applications_of_higher_computability.pdf}
\end{enumerate}

The user reports that test\_alpha0, the originator of F, circulated a separate proof
on QQ and acknowledged the independent development of G. An exact
public date and stable public link have not been established for that
communication. The acknowledgment is recorded as one of independence,
not as verification of this manuscript. The initial idea of G is
attributed to the user; the subsequent AI-assisted arguments are not
presented as the user's unaided work.
